\chapter{}

Pietro, lontano dall’uscio, ad ogni passo che udiva sperava che fosse Ghìsola; finalmente, la vide.

Non pronunciarono nè meno una parola: c’era tra loro una specie d’ostilità rispettosa. Ella volgeva gli occhi attorno; ed egli seguitava sempre i suoi occhi che lo evitavano, quantunque paresse che lo vedessero lo stesso. Tuttavia, da poche parole che avevano dovuto dirsi, sentirono svanire il loro ritegno.

Quando il tranvai si fermò, salirono.

Ella aveva un cappello di paglia, con un solo nastro di velluto nero; una veletta chiara sul volto, i guanti di filo bianco. Pietro s’accorse di quell’eleganza grossolana; e perchè se ne sentì commosso, le toccò una mano. Egli, certo, sposatala, l’avrebbe fatta vestire molto meglio. Ma tutti la guardavano; ed egli ne era contento per lei.

Andarono in fretta dalla Piazza del Duomo alla stazione, perchè c’era poco tempo alla partenza del treno. Nelle vie la folla li faceva sovvenire di sè stessi e della loro decisione, come se trasalissero. E, allora, si guardavano negli occhi. Ma presero lo stesso il treno per Siena; quasi senza parlarsi mai. Soltanto quando il loro scompartimento fu più vuoto, egli le disse:

— Perchè non t’alzi la veletta?

E le soggiunse sottovoce:

— Ti vedrò meglio.

Ella obbedì; e si sederono l’uno di fronte all’altro.

— Se ti vuoi riposare, vengo vicino a te. Vuoi appoggiare la testa su la mia spalla?

— Non importa.

Si sentivano legati dai loro sguardi, come dalle loro anime; che parevano pesanti.

Tutta la campagna correva, correva troppo! Pareva a Pietro che lo sfuggisse e non lo volesse comprendere più; anzi, lo disapprovasse. E allora aveva più bisogno d’amare Ghìsola.

Ma il giorno veniva meno come la sua esaltazione: la mattina, nel sole chiaro, gli era parso che i vagoni fossero per bruciare e fiammeggiare; ora, gli pareva, ad ogni stazione, che avessero paura di restare negli altri binarii, tutti intrecciati, dritti e curvi; che luccicavano una triste luce morta portandola con sè nell’oscurità delle lontananze diafane. La campagna si cambiava come i suoi stati mentali; ma non gli apparteneva.

A Poggibonsi, un treno, allontanandosi, divenne a poco a poco più corto, finchè non ne restò che l’ultimo vagone visto di dietro; e non si sapeva più se stesse fermo o se camminasse: come certe sue illusioni. I vagoni che andavano su e giù, trainati, con le ruote che giravano con movimento eguale l’una dopo l’altra su le medesime rotaie, e i vagoni di un treno merci verniciati di rosso, con le cifre in bianco sigillati, pazienti, lo fecero quasi piangere. Tutti scuotevano la sua anima, la schiacciavano!

Egli si sentì proprio solo e abbandonato e non si ricordò più di Ghìsola che, seduta dinanzi, lo guardava con acuta curiosità; e allora i suoi occhi avevano una immobilità affascinante.

Quand’egli, dopo aver sospirato, glieli vide così, esclamò:

— Oggi mi vuoi più bene!

Ella lo fissò con disprezzo; ma abbassò in fretta le palpebre, per nascondere lo sguardo: se le sentiva come portare via dall’anima.

Il giovine, senza capire, attese che parlasse lei, ora.

Allora Ghìsola lo fece sedere accanto; e si tennero per mano.

La gente che saliva e scendeva dal treno, i segnali delle stazioni le aumentavano la noia.

A Siena, ricusò di andare in casa della zia.

— Ma perchè non vuoi?

— Vorrà sapere troppe cose, da me: io agli altri non voglio dir niente di me.

Ella ci riusciva a vivere come voleva! La sentiva forte e indipendente. Ma per assicurarsi che non lo faceva per nascondere qualche cosa, le disse:

— Fai male: è la tua zia.

— Se andassi ad un albergo?

— Vedendoti sola penserebbero male di te.

— E tu non sai ch’io sono tua?

E insistè con un tono di voce quasi infantile, con certe moine; battendogli il ventaglio sopra un braccio:

— Sì: accontentami. Vuoi fare sempre a modo tuo. Non è vero che questa sera accontenterai la tua Ghìsola?

Volevano decidersi, perchè la strada fino alla trattoria era corta e già faceva oscuro.

Videro, dietro la basilica di San Francesco, una sfilata bassa di nuvole come il fuoco.

Qualcuno rallentava il passo per guardarli meglio, e allora camminavano più in fretta.

Alla loro sinistra si scoprì una parte di Siena, con la chiesa della Madonna di Provenzano. Tutte le case sembravano troppo fitte.

Ambedue, senza accorgersene, smisero di parlare. La Via Vallerozzi sembrava una scalinata di tetti larghi fino all’antica rocca dei Salimbeni; il cui sprone era coperto dall’ombra nera di un abete enorme. Di là da questa rocca, non si sa dove, la cima della Torre; e, più discosto, la cupola della Madonna di Provenzano, quasi rinchiusa dentro un’altra spianata di case. Mentre i tetti delle tre vie, che s’annodano insieme a Porta Ovile, scendevano, pendendo tutti da una parte; come se le case non potessero stare dritte. Un pezzetto d’una delle vie assomigliava a un baratro pietroso; e una donna, ferma, vi sembrava rinchiusa.

Tutti quei tetti, ad angolo, s’appiattivano; e alla casa più bassa, all’ultima, s’appoggiava tutta la fila delle altre.

Pietro, interrompendo la distrazione, la scosse per una mano e riprese:

— Scusami se non voglio.... Ma dài retta a me.

Ella s’impazientì e si fermò un’altra volta.

— Ascoltami.... ho pensato di portarti a mangiare da mio padre. Io gli ho detto che andavo a Poggibonsi, dove ho un amico; e gli inventerò che ti ho trovata in treno.

Ella aspettò che uno smettesse di guardarla, e poi rispose:

— E crederà a noi?

Già la curiosità dei passanti li impacciava con molestia, con tedio penoso.

— Certamente!

Ghìsola stette molto tempo a testa bassa, non per riflettere, ma per sforzarsi a non pensare ad altro: e poi rispose:

— Mi piace poco.

Tacquero perchè si sentirono vicini a bisticciare; poi egli, dopo uno di quei silenzii in cui si odono tutte le cose, la prese a braccetto fino allo scalino della trattoria.

Domenico quando li vide entrare, salutò Ghìsola ma senza avvicinarsi; e credette lì per lì alla scusa di Pietro; che del resto non aveva mentito mai.

Il marito di Rebecca, con un piatto in mano, si fermò e le disse:

— A pena che avrò servito questi signori, avvertirò la tua zia.

Ghìsola, vedendo come la parente le potesse esser di pretesto per essere venuta a Siena, lo ringraziò.

Domenico, ch’era di buon umore, dopo averla guardata sorridendo, così irriconoscibile da quando stava a Poggio a’ Meli, andò in cucina; e, come se si fosse trattato di avventori, ordinò a voce alta da cena per Pietro e per lei. Ma disse anche per farsi intendere subito:

— Questi non pagano!

Ghìsola, disinvolta, si mise a ridere; e le dispiacque solo per orgoglio che Domenico la trattasse per quel che era; ma Pietro le fece rabbia. Non era punto furbo, e non contava proprio niente in casa sua!

Per far vedere che non aveva bisogno di mangiare in trattoria, non voleva sedersi a tavola; ma Pietro la supplicò, sottovoce, di non insistere; e le disse che il giorno dopo le cose sarebbero state chiarite.

Domenico, che veniva e andava dalla cucina alla stanza dov’erano essi, con le mani in tasca e con la testa bassa, senza guardarli mai, uscì e andò sfogarsi dal suo amico droghiere: un figliolo non doveva portarsi in casa le amanti, sia pure che facesse bene a fare il comodo suo ora che era giovine. Ma il droghiere rise della sua collera e gli disse che lo lasciasse divertire, giacchè si trattava di una bella ragazza.

Ghìsola, mangiando, non alzò mai la testa; e pareva che avesse poco appetito. Ma Pietro le pestava leggermente i piedi e le diceva qualche parola perchè dissipasse il malumore. Poi la lasciò nella trattoria a chiacchierare con la cugina Rosaura, accanto alla dispensa, dov’era meno luce. E Ghìsola accompagnata da lei andò a trovare la zia, raccontandole una filza di abili menzogne, con l'aria più ingenua che ci fosse. Rebecca le disse:

— Per stasera, non ho da darti da dormire qui. Dormirai con la tua cugina, se il padrone è contento.

Ghìsola ridiscese ed entrò nella bottega, curiosa di vedere come sarebbe andata a finire!

Già era prossima la mezzanotte; e le tavole della trattoria sparecchiate. I cuochi sonnecchiavano appoggiati al ceppo del tagliere. I fornelli si spegnevano: come se anche la brace s’addormentasse. Tutti i lumi abbassati; e la trattoria piena di quell’odore ripugnante di tante vivande insieme.

In un corbello vicino all’acquaio, le bucce delle frutta e gli avanzi.

Improvvisamente la notte si fece più oscura e piovve alcuni minuti: una di quelle piogge che fanno notare subito il nostro malumore, come quelle che ribollono l’immondizie ammucchiate in mezzo ai campi.

A Ghìsola, presa dalla stanchezza e dal sonno, parve che piovesse nella sua anima, ma non riesciva a togliere tutte le cose che c’erano. Si sentiva soffocare lo stesso.

Qualche lampo, silenzioso, s’accese tra le nuvole.

Allora, ella credette che avrebbe risentito quella pioggia in qualche sogno. Evitava di non pensarci, per essere attenta a quel che accadeva intorno a lei e a quel che le dicevano.

Domenico, svegliatosi dal canapè dove da qualche tempo dormiva almeno due ore prima d'andare a letto, ordinò:

— Chiudete le porte.

Era evidente la sua scontentezza; tanto che Rosaura non gli disse volentieri:

— Io salgo in casa a trovare le lenzuola per Ghìsola.

Domenico non disse nè sì nè no; e si volse dalla parte opposta quando Ghìsola passando rapidamente vicino a lui, quasi provocandolo, lo salutò.

La camera di Rosaura era così bassa che, stando sdraiati su uno di quei letti, si poteva toccare una trave. Una finestra strettissima, nel muro più grosso di un metro, dava in una corte angusta e umida anche d'estate.

Messe le lenzuola, Ghìsola togliendosi la giacchetta domandò:

— Dove dorme Pietro?

— Nella stanza solita di quando era piccino. Ma vorresti andare a vederlo? Che braccia grosse tu hai!

— Senti come sono ingrassata!

Si fece pizzicare un fianco, e poi andò.

Riconoscendo bene la casa, si avanzò quasi a tastoni, attraversando la stanza d'ingresso e poi il salotto meno buio perchè c'era la luce elettrica dalla strada.

L'uscio della camera di Pietro era aperto perchè vi doveva passare Domenico per andare nella sua. Ella vide il tavolino con i libri, il canterano con lo specchio che luccicava. E proseguì verso il tetto messo ad una parete: Pietro dormiva.

Allora si chinò e cominciò a baciarlo su la bocca. Egli, senza finire di destarsi, sentì un brivido; ed esclamò a voce alta:

— Sei tu, Ghìsola!
